\begin{table}[!htb]
\centering
\caption{Evaluasi Alternatif \textit{Runtime Security} \textit{Open-Source} pada Kubernetes}
\label{tab:runtime_security_comparison}
\footnotesize
\setlength{\tabcolsep}{4pt}
\begin{tabular}{|>{\raggedright\arraybackslash}m{3.5cm}|>{\centering\arraybackslash}m{1.5cm}|>{\centering\arraybackslash}m{2cm}|>{\centering\arraybackslash}m{2cm}|>{\centering\arraybackslash}m{1.7cm}|>{\centering\arraybackslash}m{1cm}|}
\hline
\multicolumn{1}{|>{\centering\arraybackslash}m{3.5cm}|}{\textbf{\textit{Runtime Security}}} & \textbf{Maturitas} & \textbf{Cakupan eBPF/Kernel} & \textbf{Aturan Bawaan} & \textbf{Komunitas} & \textbf{Total} \\
\hline
Falco (CNCF) & 2 & 2 & 2 & 2 & 8 \\
\hline
Cilium Tetragon & 1 & 2 & 1 & 1 & 5 \\
\hline
Aqua Tracee & 1 & 2 & 1 & 1 & 5 \\
\hline
Sysdig Secure \textit{OSS} & 1 & 1 & 1 & 1 & 4 \\
\hline
\end{tabular}
\end{table}
